% concatenated from samplepaper.tex
% This is samplepaper.tex, a sample chapter demonstrating the
% LLNCS macro package for Springer Computer Science proceedings;
% Version 2.21 of 2022/01/12
%
\documentclass[runningheads]{llncs}
%
\usepackage[T1]{fontenc}
% T1 fonts will be used to generate the final print and online PDFs,
% so please use T1 fonts in your manuscript whenever possible.
% Other font encondings may result in incorrect characters.
%
\usepackage{graphicx}
% Used for displaying a sample figure. If possible, figure files should
% be included in EPS format.
%
% If you use the hyperref package, please uncomment the following two lines
% to display URLs in blue roman font according to Springer's eBook style:
%\usepackage{color}
%\renewcommand\UrlFont{\color{blue}\rmfamily}
%
\usepackage{amsmath}
\usepackage{amssymb}
\usepackage{mathtools}
\usepackage{enumitem}

\spnewtheorem{notation}{Notation}{\bfseries}{}

\newcommand{\N}{\mathbb{N}}
\newcommand{\Z}{\mathbb{Z}}
\newcommand{\R}{\mathcal{R}}
\newcommand{\set}[2]{\{ #1 \: | \: #2 \}}
\newcommand{\confspace}{\Sigma^{\Z^d}}
\newcommand{\id}{\mathrm{id}}
\newcommand{\cyl}{\mathrm{Cyl}}

\begin{document}
%
\title{Reversibility, balance and expansivity of non-uniform cellular automata\thanks{Supported by the Magnus Ehrnrooth foundation and the Research Council of Finland.}}
%
%\titlerunning{Abbreviated paper title}
% If the paper title is too long for the running head, you can set
% an abbreviated paper title here
%
\author{Katariina Paturi\orcidID{0009-0007-2475-1393}}
%
\authorrunning{K. Paturi.}
% First names are abbreviated in the running head.
% If there are more than two authors, 'et al.' is used.
%
\institute{University of Turku, Vesilinnantie 5, 20500 Turku, Finland\\
\email{katariina.paturi@utu.fi}
}
%
\maketitle              % typeset the header of the contribution
%
\begin{abstract}
Non-uniform cellular automata (\textsc{nuca}) are an extension of cellular automata (\textsc{ca}), which transform cells according to multiple different local rules.  A \textsc{nuca} is defined by a configuration of local rules called a local rule distribution. We examine what properties of uniform \textsc{ca} can be recovered by restricting the rule distribution to be (uniformly) recurrent, focusing on only 1D \textsc{nuca}.

We show that a bijective \textsc{nuca} with a uniformly recurrent rule distribution is reversible. We also show that if a \textsc{nuca} is surjective and has a recurrent rule distribution, or if it is bijective, then it is balanced. We present an example of a \textsc{nuca} which has a non-empty and non-residual set of equicontinuity points, and one which is not sensitive but has no equicontinuity points. Finally, we show that (positively) expansive \textsc{nuca} are sensitive.

\keywords{Non-uniform cellular automata \and Symbolic dynamics \and Reversibility.}
\end{abstract}
%
%
%
\section{Introduction}

Non-uniform cellular automata (\textsc{nuca}) are an extension of cellular automata (\textsc{ca}), which are machines that operate on configurations - grids of cells coloured by an alphabet. In both, each cell simultaneously observes its neighbouring cells and changes its state depending on the state of the neighbourhood. A uniform \textsc{ca} transforms all cells with the same local rule, while while a \textsc{nuca} uses different local rules in different cells. The local rules of a \textsc{nuca} are given by a fixed local rule distribution. 

Cellular automata are a powerful tool in simulating systems based on local transformations, and \textsc{nuca} provide a generalization capable of simulating systems with variation in their local behaviour. The dynamical properties of uniform \textsc{ca} have been studied extensively, but many of them no-longer apply in the setting of \textsc{nuca}. For example, the well-known Garden of Eden theorem for uniform \textsc{ca} does not hold for \textsc{nuca}.  However, some properties, such as the Garden of Eden theorem, can be recovered in \textsc{nuca} when the distribution of local rules is recurrent or uniformly recurrent \cite{goe}. 

The properties we examine are reversibility, balance, sensitivity and expansivity. There are implications known to hold for \textsc{ca} which don't necessarily hold for \textsc{nuca} in general: all bijective \textsc{ca} are reversible, all surjective \textsc{ca} are balanced, a \textsc{ca} is either sensitive  or has a dense set of equicontinuity points, and a (positively) expansive \textsc{ca} is sensitive. We find (positively) expansive \textsc{nuca} are still sensitive, but the other statements all have counter-examples. We also show that the reversibility of bijective \textsc{nuca} can be recovered when the rule distribution is uniformly recurrent, and the balance of surjective \textsc{nuca} when the distribution is recurrent. Additionally we show that bijective \textsc{nuca} are balanced. 

\section{Definitions}

Non-uniform cellular automata may be defined over any group, but in this work we only consider the 1-dimensional case of \textsc{nuca} over $\Z$. First, a few background definitions are needed.

\begin{definition}
Let $\Sigma$ be a finite set called a \emph{state set}. A \emph{configuration (on $\Z$)}  is an element $c \in \Sigma^\Z$.
\end{definition}

\begin{definition}
A \emph{cell} is an element $x\in \Z$. A \emph{domain} is any subset $D\subseteq \Z$. If $D$ is finite, we call it a \emph{finite domain}. 

Let $\Sigma$ be a state set and $c \in \Sigma^\Z$ a configuration. A \emph{pattern} is any $p \in \Sigma^D$ for some domain $D$. If $D$ is finite, we call $p$ a \emph{finite pattern}. We use $c_{|D}\in \Sigma^D$ to denote the pattern such that for any $x\in D$, $c_{|D} (x) = c(x)$.
\end{definition}

%\begin{definition}
%Let $k\in \Z_+$. A \emph{neighbourhood} is a tuple $N = (x_1,\ldots,x_k)\in \Z^k$. If $r\in \Z_+$, a \emph{radius-$r$ neighbourhood} is the neighbourhood $N_r = (-r,-r+1,\ldots,0,\ldots,r-1,r)$.  For a cell $x \in \Z$, we denote
%\begin{align*}
%N(x) = \set{x+x_i}{x_i \in N},
%\end{align*}
%and call the domain $N(x)$ the \emph{neighbourhood of $x$}. Similarly, if $D\subseteq \Z$ is a domain, we denote
%\begin{align*}
%N(D) = \set{x+x_i}{x\in D, x_i \in N},
%\end{align*}
%and call the domain $N(D)$ the \emph{neighbourhood of $D$}. 
%\end{definition}

The space $\Sigma^\Z$ equipped with Cantor's topology and is known to be compact \cite{cecc}. In contrast with uniform \textsc{ca}, which are defined by a local rule, a \textsc{nuca} is defined by a local rule distribution, which is a special configuration that determines which local rule is used in which cell.

\begin{definition}
Let $\Sigma$ be a state set. A \emph{radius $r$ local rule} is a function $f:\Sigma^{2r+1} \rightarrow \Sigma$. A finite set of local rules $\mathcal{\R}$ is called a \emph{rule set}.
\end{definition}

Note that defining local rules with radial neighbourhoods is equivalent to defining them with arbitrarily shaped neighbourhoods. Now we are ready to define \textsc{nuca}.

\begin{definition}
A \emph{non-uniform cellular automaton} (or \textsc{nuca}) is a tuple $A = (\Sigma, \R,\theta)$, where $\Sigma$ is a state set, $\R$ is a rule set, and $\theta \in \R^\Z$ is a rule distribution. The \emph{global transition function} of $A$ is the function $H_\theta:\Sigma^\Z \rightarrow \Sigma^\Z$ such that for $c\in \Sigma^\Z$, $x \in \Z$, 
\begin{align*}
H_\theta (c)(x) = \theta(x) (c(x-r),\ldots,c(x+r)),
\end{align*}
where $r \in \Z_+$ is the radius of the local rule $\theta (x)$. If $|\R| = 1$ then the \textsc{nuca} is a uniform cellular automaton (or \textsc{ca}).
\end{definition}

We generally equate a \textsc{nuca} with its transition function. Let us define a few more concepts used throughout the work. The definitions of recurrence and uniform recurrence are equivalent with the usual topological definitions for subshifts.

\begin{definition}
Let $\theta \in \R^\Z$ be a rule distribution, and suppose $\theta (x)$ is a radius $r$ rule. We call the domain $N_\theta(x)=[x-r,x+r]$ the \emph{neighbourhood of $x$}. If $D \subseteq \Z$ is a domain, we call the domain $N_\theta (D) = \bigcup_{x\in D} N_\theta (x)$ the \emph{neighbourhood of $D$}.
\end{definition}

\begin{definition}
Let $\Sigma$ be a state set. The \emph{left-shift} is the function $\sigma:\Sigma^\Z \rightarrow \Sigma^\Z$ which maps $c \in \Sigma^\Z$ such that $\sigma (c)(x) = c(x+1)$ for all $x\in \Z$. Similarly, we call the function $\sigma^{-1}$ the \emph{right shift}. Note that $\sigma \circ H_\theta = H_{\sigma(\theta)} \circ \sigma$.

If $p\in \Sigma^D$ and $q\in \Sigma^C$ are patterns and there is $k \in \Z$ such that $C = D+k$ and $p(x) = q(x+k)$ for any $x \in D$, then we call $q$ a \emph{translated copy} (or simply \emph{copy}) of $p$. If $D \cap (D+k) = \emptyset$, the copies are \emph{disjoint}.
\end{definition}

\begin{definition}
Let $c \in \Sigma^\Z$ be a configuration. The configuration $c$ is \emph{(spatially) recurrent} if for any finite $D\subseteq \Z$ there is $C \subseteq \Z$ such that $D \neq C$ and $c_{|C}$ is a copy of $c_{|D}$. The configuration $c$ is \emph{(spatially) uniformly recurrent} if for any finite $D\subseteq \Z$ there is $n \in \Z_+$ such that for any $x \in \Z$, there is $C \subseteq [x,x+n]$ such that  $c_{|C}$ is a copy of $c_{|D}$. The same concepts are defined for rules rule distributions $\theta \in  \R^\Z$ when $\R$ is finite.
\end{definition}

\section{Reversibility}

Informally, we call a \textsc{nuca} $H$ reversible if there exists an inverse \textsc{nuca} $H^{-1}$, such that $H\circ H^{-1} = H^{-1} \circ H$ is the identity function. In uniform \textsc{ca} reversibility is equivalent with bijectivity, which in turn is equivalent with injectivity. 

\begin{definition}
Let $\Sigma$ be a state set. By $\id:\Sigma^\Z \rightarrow \Sigma^\Z$ we denote the identity function.
\end{definition}

\begin{definition}
Let $\R_1$ be a rule set and $\theta \in \R_1^\Z$. If there exists finite rule set $\R_2$ and $\phi \in \R_2^\Z$ such that $H_\phi = H_\theta^{-1}$, then the automaton $H_\theta$ is \emph{reversible}.
\end{definition}

Obviously, a reversible \textsc{nuca} is bijective. For uniform \textsc{ca}, the converse requires a non-trivial proof based on the compactness of the configuration space $\Sigma^\Z$. If we allow \textsc{nuca} to have infinitely many rules, the converse is also simple. This follows from the fact that \textsc{nuca} (allowing infinitely many rules) are exactly the continuous functions over the Cantor set \cite{dennuz}, and the fact that the inverse of a bijective continuous function is also continuous. Then the inverse of a bijective \textsc{nuca} is a continuous function over the Cantor set, thus also being a \textsc{nuca}.

%\begin{theorem} \label{inf_revers}
%Let $\R$ be a (potentially infinite) ruleset and $\theta \in \R^\Z$ such that $H_\theta$ is bijective. Then $H_\theta$ is reversible.
%\end{theorem}

It is then natural to ask if this is true restricting ourselves to \textsc{nuca} with finitely many rules only. The following is an example of a bijective \textsc{nuca} which is not reversible. 


\begin{example}
We denote $a \oplus b = a + b \mod 2$. Let $\Sigma = \Z_2$ and $\R = \{f_L, f_R, g \}$ be a set of radius $1$ rules over $\Sigma$, which map
\begin{align*}
f_L(a,b,c) &= a \oplus b, \\
f_R(a,b,c) &= b \oplus c, \\
g(a,b,c) &= b,
\end{align*}
where $a,b,c \in \Sigma$. For all $n \in \Z_+$, let $u_n = (f_R)^n g (f_L)^n$ and let $w_n =  u_1 \ldots u_n$. Then, let $\theta = \ldots w_3 w_2 w_1 w_2 w_3 \ldots$ be the rule distribution.

The distribution $\theta$ is clearly recurrent. First, let's show that $H_\theta$ is injective. Let $c,e \in \Sigma^\Z$ be such that $c(x) \neq e(x)$ for some $x \in \Z$. If $\theta (x) = g$, then $H_\theta (c)(x) = c(x)\neq e(x)= H_\theta (e)(x)$. Suppose then $\theta (x) = f_R$. Then there is $k\in \Z_+$ such that  $\theta(x+i) = f_R$ when $0 \leq i < k$ and $\theta(x+k) = g$. 

Now suppose $H_\theta (c)(x+i) = H_\theta (e)(x+i)$ when $0 \leq i < k$. If $c(x+j) \neq e(x+j)$ where $0 \leq j < k$, then $c(x+j+1) \neq e(x+j+1)$ also, because otherwise it would hold that
\begin{align*}
H_\theta (c)(x+j) &= c(x+j) \oplus c(x+j+1) \\
&= c(x+j) \oplus e(x+j+1)\\
&\neq e(x+j) \oplus e(x+j+1) = H_\theta(e)(x+j)
\end{align*}
Now inductively, $c(x+i) \neq e(x+i)$ when  $0 \leq i \leq k$. Then because $c(x+k) \neq e(x+k)$ and $\theta(x+k)=g$, we have $H_\theta(c)(x+k) \neq H_\theta(e)(x+k)$. In any case, $H_\theta (c) \neq H_\theta (e)$.

The case that $\theta(x) = f_L$ is symmetrical. Therefore $H_\theta$ is injective. Because $\theta$ is recurrent, by the Garden of Eden theorem, $H_\theta$ is then also bijective \cite{goe}. Next we show that $H_\theta$ is not reversible with a finite rule set. 

Suppose there is a finite rule set $\R'$ and $\phi \in \R'^\Z$ such that $H_\phi = H_\theta^{-1}$. Then there is some $r\in \Z_+$ such that any rule in $\mathcal{R}'$ is at most radius $r$. Somewhere in $\theta$, there is a length  $2r+1$ segment containing only the rule $f_R$; let $x \in \Z$ be the centermost cell of such a segment. Let then $c,e \in \Sigma^\Z$ be such that $c(y) = 0$ and $e(y) = 1$ for all $y\in \Z$. Then for all $z \in [x-r,x+r]$,
\begin{align*}
H_\theta(c)(z) = 0 \oplus 0 = 1 \oplus 1 = H_\theta(e)(z).
\end{align*}
Suppose $\phi(x)(0,\ldots,0) = a$ for some $a \in \Sigma$. Then 
\begin{align*}
c(x) = (H_\phi \circ H_\theta) (c) (x) = a = (H_\phi \circ H_\theta) (e) (x) = e(x),
\end{align*}
which is a contradiction. Therefore $H_\theta$ is not finitely reversible.
\end{example}

Note that the rule distribution in the previous counter-example is recurrent. A simpler, yet non-recurrent counter-example in the notation of the previous example would be distribution $^\infty (f_R)g(f_L)^\infty$. Next we show that if a bijective \textsc{nuca} is defined by a uniformly recurrent rule distribution, then it is reversible. To do so, we use the notion of stable injectivity. 

\begin{definition}
Let $\theta\in \R^\Z$ be a rule distribution. The \emph{orbit} of $\theta$ is the set $\mathcal{O}(\theta) = \set{\sigma^k (\theta)}{k\in \Z}$. The \emph{orbit closure} $\overline{\mathcal{O}(\theta)}$ is the topological closure of $\mathcal{O}(\theta)$. We call $H_\theta$ \emph{stably injective} if for any $\phi \in \overline{\mathcal{O}(\theta)}$, the \textsc{nuca} $H_\phi$ is injective. \cite{phung}
\end{definition}

\begin{lemma} \label{lemma_inj}
Let $\theta \in \R^\Z$ be a uniformly recurrent rule distribution such that $H_\theta$ is injective. Then $H_\theta$ is stably injective.
\end{lemma}

\begin{proof}
Let $\phi \in \overline{\mathcal{O}(\theta)}$. Suppose $H_\phi$ is not injective. Then there are two cases: suppose first that $H_{\phi}$ is not pre-injective. Let then $c'',e'' \in \Sigma^\Z$ be asymptotic configurations such that $c''\neq e''$ and $H_{\phi} (c'') = H_{\phi} (e'')$, and $x,y\in \Z$ such that for any cell $z \notin [x,y]$, $c''(z) = e''(z)$. Because $\phi \in \overline{\mathcal{O}(\theta)}$, any pattern in $\phi$ has a copy in $\theta$. Denote $D=[x-r,y+r]$ and let $m \in \Z$ be such that $\theta_{|(D+m)}$ is a copy of $\phi_{|D}$. Then because for any cell $z \in (D+m)$, $\theta(z)=\phi(z-m)$, so 
\begin{align*}
H_\theta (\sigma^{-m} (c''))(z) = H_{\phi} (c'')(z-m) = H_{\phi} (e'')(z-m) =H_\theta (\sigma^{-m} (e''))(z),
\end{align*}
and for any cell $z \notin (D+m)$, $\sigma^{-m} (c'')_{|[z-r,z+r]}=\sigma^{-m} (e'')_{|[z-r,z+r]}$, we have $H_\theta (\sigma^{-m} (c'')) = H_\theta (\sigma^{-m} (e''))$. Then $H_\theta$ is not pre-injective, which is a contradiction.

Suppose then that $H_{\phi}$ is pre-injective. Then either every length $2r+1$ segment to the right of the origin 0 contains a cell where $c'$ and $e'$ differ, or the same holds for every $2r+1$ segment to the left of the origin. Assume without loss of generality the first case, that is for every $k>0$ there is $i \in [-r,r]$ such that $c'(k+i) \neq e'(k+i)$.

Because $\theta$ is uniformly recurrent, so is $\phi$ and $\theta \in \overline{\mathcal{O} (\phi)}$. Hence every finite segment of $\theta$ appears in $\phi$ in every segment of some length $n>0$. Then there exists $k_1, k_2, k_3, \ldots > 0$ such that $\sigma^{k_1} (\phi),\sigma^{k_2} (\phi),\ldots $ converges to $\theta$. Select then a subsequence  $k_{i_1},k_{i_2},\ldots$ such that $c'^{k_{i_1}},c'^{k_{i_2}},\ldots$ and $e'^{k_{i_1}},e'^{k_{i_2}},\ldots$ converge to some $c''$ and $e''$ respectively. Then by the previous we have $c''_{|[-r,r]}\neq e''_{|[-r,r]}$, and by continuity of \textsc{nuca}, $H_\theta (c'') = H_\theta (e'')$. Then $H_\theta$ is not injective, which is a contradiction. 
\qed
\end{proof}

Note that Lemma \ref{lemma_inj} doesn't hold in general, even in $\Z^2$. Since bijective, stably injective NUCA are reversible \cite{phung}, we can prove the following theorem.

\begin{theorem} \label{unif_revers_fin}
Let $\R$ be a finite rule set and $\theta\in \R^\Z$ a uniformly recurrent rule distribution such that $H_\theta$ is injective. Then $H_\theta$ is reversible.
\end{theorem}

\begin{proof}
 By Lemma \ref{lemma_inj}, $H_\theta$ is stably injective. Because $\theta$ is recurrent and $H_\theta$ is injective, then $H_\theta$ is bijective \cite{goe}. Then there is a finite rule set $\mathcal{R}_2$ and $\phi \in \R_2^\Z$, such that $H_\phi \circ H_\theta = \id$ \cite{phung}. Now because $H_\theta$ is surjective,  we have $H_\phi = H_\theta^{-1}$.

%We may assume all rules of $\R_1$ have radius $r$. Suppose to the contrary that for any $R \in \N$ there is $k \in \Z$ and $c,e \in \Sigma^\Z$ such that $c(k) \neq e(k)$ and $H_\theta (c)_{|[k-R,k+R]}=H_\theta (e)_{|[k-R,k+R]}$, in other words, a radius $R$ window is not enough to differentiate the images of some configurations at cell $k$.
%
%By shifting cell $k$ to the origin 0 for each $R$, we obtain sequences $c_1,c_2,\ldots$, $e_1,e_2,\ldots$, $\theta_1, \theta_2,\ldots$ such that for any $i \in \Z_+$,
%\begin{align*}
%c_i(0) &\neq e_i(0), \\
%H_{\theta_i} (c_i)_{|[-i,i]} &= H_{\theta_i} (e_i)_{|[-i,i]},\\
%\theta_i &\in \mathcal{O}(\theta).
%\end{align*}
%By compactness the sequences have simultaneously converging sub-sequences with respective limits $c',e', \theta'$, which then satisfy
%\begin{align*}
%c'(0) &\neq e'(0), \\
%H_{\theta'} (c') &= H_{\theta'} (e'),\\
%\theta' &\in \overline{\mathcal{O}(\theta)},
%\end{align*}
%meaning $H_{\theta'}$ is not injective. This contradicts Lemma \ref{lemma_inj}, and hence the theorem statement holds.
%
%Then there are two cases: suppose first that $H_{\theta'}$ is not pre-injective. Let then $c'',e'' \in \Sigma^\Z$ be asymptotic configurations such that $c''\neq e''$ and $H_{\theta'} (c'') = H_{\theta'} (e'')$, and $x,y\in \Z$ such that for any cell $z \notin [x,y]$, $c''(z) = e''(z)$. Because $\theta' \in \overline{\mathcal{O}(\theta)}$, any pattern in $\theta'$ has a copy in $\theta$. Denote $D=[x-r,y+r]$ and let $m \in \Z$ be such that $\theta_{|(D+m)}$ is a copy of $\theta'_{|D}$. Then because for any cell $z \in (D+m)$, $\theta(z)=\theta'(z-m)$, so 
%\begin{align*}
%H_\theta (\sigma^{-m} (c''))(z) = H_{\theta'} (c'')(z-m) = H_{\theta'} (e'')(z-m) =H_\theta (\sigma^{-m} (e''))(z),
%\end{align*}
%and for any cell $z \notin (D+m)$, $\sigma^{-m} (c'')_{|[z-r,z+r]}=\sigma^{-m} (e'')_{|[z-r,z+r]}$, we have $H_\theta (\sigma^{-m} (c'')) = H_\theta (\sigma^{-m} (e''))$. Then $H_\theta$ is not pre-injective, which is a contradiction.
%
%Suppose then that $H_{\theta'}$ is pre-injective. Then either every length $2r+1$ segment to the right of the origin 0 contains a cell where $c'$ and $e'$ differ, or every the same holds for every $2r+1$ segment to the left of the origin. Assume without loss of generality the first case, that is for every $k>0$ there is $i \in [-r,r]$ such that $c'(k+i) \neq e'(k+i)$.
%
%Because $\theta$ is uniformly recurrent, then $\theta \in \overline{\mathcal{O} (\theta')}$ and hence every finite segment of $\theta$ appears in $\theta'$ in every segment of some length $n>0$. Then there exists $k_1, k_2, k_3, \ldots > 0$ such that $\sigma^{k_1} (\theta'),\sigma^{k_2} (\theta'),\ldots $ converges to $\theta$. Select then a subsequence  $k_{i_1},k_{i_2},\ldots$ such that $c'^{k_{i_1}},c'^{k_{i_2}},\ldots$ and $e'^{k_{i_1}},e'^{k_{i_2}},\ldots$ converge to some $c''$ and $e''$ respectively. Then by the previous we have $c''_{|[-r,r]}\neq e''_{|[-r,r]}$, and by continuity of \textsc{nuca}, $H_\theta (c'') = H_\theta (e'')$. Then $H_\theta$ is not injective, which is a contradiction. Hence the theorem statement holds.
\qed
\end{proof}
 
We can also show that the inverse automaton of a reversible \textsc{nuca} with a uniformly recurrent rule distribution will also have a uniformly recurrent rules distribution, up to renaming rules.

\begin{definition}
Let $f:\Sigma^{2r+1} \rightarrow \Sigma$ and $g:\Sigma^{2R+1} \rightarrow \Sigma$ be radius $r$ and $R$ local rules respectively for some $r,R\in \Sigma^\Z$. Assume $r \leq R$. The rules $f$ and $g$ are considered \emph{identical}, if for any $a_{-R},\ldots,a_{R} \in \Sigma$ it holds that
\begin{align*}
g(a_{-R},\ldots,a_{R}) = f(a_{-r},\ldots,a_r).
\end{align*}
\end{definition}



\begin{theorem} \label{unif_revers_unif}
Let $\R_1$ be a finite rule set and $\theta \in \R_1^\Z$ uniformly recurrent. Let $\R_2$ be a rule set containing no pair of different identical rules, and $\phi \in \R_2^\Z$ a distribution such that $H_\phi = H_\theta^{-1}$. Then $\phi$ is uniformly recurrent.
\end{theorem}

\begin{proof} 
Let $r\in \Z_+$ be at least as large as the radius of any rule in $\R_1$ or $\R_2$. We show that for any $x,y \in \Z$, if $\theta_{|[x-r,x+r]} = \sigma^k (\theta_{|[y-r,y+r]})$, where $k=y-x$, then $\phi(x) = \phi(y)$. For all $c \in \Sigma^\Z$ we have $H_\theta (c)_{[x-r,x+r]} = H_{\sigma^k (\theta)} (c)_{[x-r,x+r]}$. Then applying the local rule $\phi(x)$, we have
\begin{align*}
H_\phi (H_\theta (c)) (x) = H_\phi (H_{\sigma^k (\theta)} (c)) (x).
\end{align*}
Since $H_\phi \circ H_\theta = \id = H_{\sigma^k (\phi)} \circ H_{\sigma^k (\theta)}$, we then have
\begin{align*}
H_{\sigma^k(\phi)} (H_{\sigma^k (\theta)} (c)) (x) = H_\phi (H_{\sigma^k (\theta)} (c)) (x).
\end{align*}
Because $H_{\sigma^k (\theta)}$ is surjective, for any $e \in \Sigma^\Z$ there is $c \in \Sigma^\Z$ such that $H_{\sigma^k (\theta)}(c) = e$. Therefore for any $e\in \Sigma^\Z$, 
\begin{align*}
H_{\sigma^k(\phi)} (e) (x) = H_\phi (e) (x).
\end{align*}
Since there are no identical rules in $\R_2$, $\sigma^k(\phi)(x) = \phi(x)$. Now because $\theta$ is uniformly recurrent, for any $m \in \Z_+$, and any segment $D \subseteq \Z$ of length $|D|=m+2r$, there is $n \in \Z_+$ such that  a copy of $\theta_{|D}$ appears in every length $n$ segment of $\theta$. Then by the previous, for every segment $D'$ of length $|D'|=m$ there is $n\in \Z_+$ such that a copy of $\phi_{|D'}$ appears in every segment of length $n$. Hence $\phi$ is uniformly recurrent.
\qed

%Let $x,y \in \Z$, $x\leq y$ and $r\in \Z_+$ at least as large as the radius of any rule in $\R_1$ or $\R_2$. Because $\theta$ is uniformly recurrent, there is $n \in \Z_+$ such that a disjoint copy of the pattern $\theta_{|[x-2r,y+2r]}$ appears in every segment of length $n$ in $\theta$. Let $k\in \Z$ be such that $\theta_{|[x-2r+k,y+2r+k]}$ is a disjoint copy of $\theta_{|[x-2r,y+2r]}$. We show that $\phi_{|[x+k,y+k]}$ is a copy of $\phi_{|[x,y]}$.



%Let then $e \in \Sigma^\Z$ be such that $e(z) = e(z+k)$ and for all $z' \in [z-r,z+r]$, $H_\theta (e)(z') = H_\theta(e)(z'+k) = c(z')$. Such a configuration exists, because $H_\theta$ is bijective and the neighbourhood $N_\theta([z-r,z+r])$ doesn't overlap with the neighbourhood  $N_\theta([z-r+k,z+r+k])$. But now
%\begin{align*}
%H_\phi(e)(z) = H_\phi(c)(z) \neq H_\phi(\sigma^{-k}(c))(z+k) = H_\phi (e)(z+k),
%\end{align*}
%meaning $H_\phi \circ H_\theta (e) \neq e$, which is a contradiction. 

%Let $z \in [x,y]$, $c\in \Sigma^\Z$ and $e = H_\phi (c)$. Because $\theta_{|[x-2r+k,y+2r+k]}$ is a copy of $\theta_{|[x-2r,y+2r]}$, then 
%\begin{align*}
%H_\theta (\sigma^{-k}(e))_{|[x-r+k,y+r+k]} = \sigma^{-k}(c)_{|[x-r+k,y+r+k]}
%\end{align*}
%Now because $H_\phi = H_\theta^{-1}$, 
%\begin{align*}
%H_\phi (\sigma^{-k}(c)) (z+k) =  \sigma^{-k}(e)(z+k) = e(z) = H_\phi (c)(z).
%\end{align*}
%Therefore there's no pattern that $\phi(z)$ and $\phi(z+k)$ map differently, and hence they are identical. Because $\R_2$ contained no pairs of different identical rules, then $\phi(z) = \phi(z+k)$. Then there is $n \in \Z_+$ such that a copy of pattern $\phi_{|[x,y]}$ appears in every segment of length $n$ in $\phi$, and $\phi$ is uniformly recurrent.
\end{proof}

%\begin{lemma} \label{unif_revers2}
%Let $\R_1$ be a finite ruleset, $\R_2$ a (potentially infinite) ruleset, both over the state set $\Sigma$, and $\theta \in \R_1^\Z, \phi \in \R_2^\Z$ uniformly recurrent distributions such that $H_\phi \circ H_\theta = \id$. Then there is $R \in \Z_+$ such  that for any $x \in \Z$ and $c,e \in \Sigma^\Z$ if $c_{|[x-R,x+R]} = e_{|[x-R,x+R]}$, then $H_\phi (c) (x) = H_\phi (e) (x)$.
%\end{lemma}
%
%\begin{proof}
%Let $r_1\in \Z_+$ be the radius of the rules in $\R_1$ and let $p$ be a length $2r_1$ segment of $\phi$. Because $\phi$ is uniformly recurrent, there is $n$ such that a copy of $p$ appears in every segment of length $n$. Let $r_2$ be the largest radius of any rule in $p$.
%
%Let $x \in \Z$. Consider the nearest disjoint copies of $p$ to the left and right of $x$ in $\phi$. Let $i \in \Z$ be the leftmost cell of the left copy and $j \in \Z$ the rightmost cell of the right copy. Then the length of the segment $|[i-r_2,j+r_2]| < 4r_1 + 2r_2 + 2n$, because the gap between the two copies of $p$ must be less than $2n$, because otherwise there would be a length $n$ segment of $\theta$ with no copy of $p$. Note that this upper bound is fixed.
%
%Suppose there are configurations $c,e \in \Sigma^\Z$ such that $c_{|[i-r_2,j+r_2]} = e_{|[i-r_2,j+r_2]}$, but $H_\phi (c)(x) \neq H_\phi (e)(x)$. The $H_\phi$-images of $c$ and $e$ in the segments on the copies of $p$ are uniquely defined by the cells at most $r_2$ away from them in $c$ and $e$ respectively. More formally, since 
%\begin{align*}
%c_{|[i-r_2,i+2r_1+r_2-1]} &= e_{|[i-r_2,i+2r_1+r_2-1]},\\
%c_{|[j-r_2-2r+1,j+r_2]} &= e_{|[j-r_2-2r+1,j+r_2]},
%\end{align*}
%and because $H_\phi$ is bijective, this means $H_\phi (c)_{|[i,i+2r_1-1]} = H_\phi (e)_{|[i,i+2r_1-1]}$ and $H_\phi (c)_{|[j-2r_1+1,j]} = H_\phi (e)_{|[j-2r_1+1,j]}$. This is illustrated in Figure \ref{fig_revers1}.
%
%\begin{figure}
%\begin{center}
%	\includegraphics[scale=0.9]{revers1v2.pdf}
%	\caption{A segment bounded by copies of the pattern $p$ on its left and right. Cells within $r_2$ of a copy of $p$ are highlighted, and they uniquely determine the $H_\phi$-image of a configuration in the cells on top of said copy of $p$.} \label{fig_revers1}
%\end{center}
%\end{figure}
%
%Let then $a \in \Sigma^\Z$ be such that
%\begin{align*}
%a (x) = 
%\begin{cases}
%H_\phi (e), \: \text{if} \: i+r_1 < x < j-r_1, \\
%H_\phi (c), \: \text{otherwise}.
%\end{cases}
%\end{align*}
%Now for all $y \in [i+r_1+1, j-r_1-1]$,
%\begin{align*}
%H_\theta (a)(y) = H_\theta \circ H_\phi (e)(y) = e(y) = c(y)
%\end{align*}
%because $N_\theta(y) \subseteq [i,j]$, and $c_{[i,j]} = e_{[i,j]}$. On the other hand, for all $y \notin [i+r_1+1, j-r_1-1]$, 
%\begin{align*}
%H_\theta (a)(y) = H_\theta \circ H_\phi (c)(y) = c(y),
%\end{align*}
%because for any such cell, $a_{|[y-r_1,y+r_1]} = c_{|[y-r_1,y+r_1]}$. Then we have that $H_\theta (a) = c$, but $a(x) = e(x) \neq c(x)$, which contradicts $H_\theta$ being bijective. This is illustrated in Figure \ref{fig_revers2}.
%
%
%\begin{figure}
%\begin{center}
%	\includegraphics[scale=0.9]{revers2.pdf}
%	\caption{The two different configurations $c,e$ and their $H_\phi$-images. The configuration $a$ and its $H_\theta$ image is shown with a highlighted line.} \label{fig_revers2}
%\end{center}
%\end{figure}
%
%Now let $R= 2r_1 + r_2 + n$. By the preceding, if $c,e \in \Sigma^\Z$ are such that $c_{|[x-R,x+R]} = e_{|[x-R,x+R]}$, then $H_\phi (c)(x) = H_\phi (e)(x)$.
%\qed
%\end{proof}
%
%By combining these lemmas with Theorem \ref{inf_revers}, we get our result about the reversibility of uniformly recurrent distributions.

%\begin{theorem}
%Let $\R_1$ be a finite rule set and $\theta \subseteq \R_1^\Z$ a uniformly recurrent rule distribution. If $H_\theta$ is bijective, there is finite $\R_2$ and $\phi \in \R_2^\Z$ such that $H_\phi \circ H_\theta = \id$.
%\end{theorem}
%
%\begin{proof}
%By Theorem \ref{inf_revers}, if $H_\theta$ is bijective, there is a potentially infinite rule set $\R_2$ and distribution $\phi \in \R_2^\Z$ such that $H_\phi \circ H_\theta = \id$. Without loss of generality, we can assume $\R_2$ has no pair of identical rules. Then by Lemma \ref{unif_revers1}, the rule distribution $\phi$ is uniformly recurrent. 
%
%Now by Lemma \ref{unif_revers2}, there is $R \in \Z_+$ such  that for any $x \in \Z$ and $c,e \in \Sigma^\Z$ if $c_{|[x-R,x+R]} = e_{|[x-R,x+R]}$, then $H_\phi (c) (x) = H_\phi (e) (x)$. Therefore every rule in $\R_2$ is identical to some radius $R$ rule. Because there is no pair of identical rules in $\R_2$, and there are only finitely many rules with radius $R$ over a fixed state set, $\R_2$ must then be finite.
%\qed
%\end{proof}


\section{Balance}

Balance is a property found in surjective \textsc{ca}, where every finite pattern in a given domain has the same number of pre-images. Balance is equivalent to invariance of the uniform measure. Surjective \textsc{nuca} are not necessarily balanced, as shown by the following example.

\begin{definition}
Let $\theta \in \R^{\Z^d}$ be a rule distribution with state set $\Sigma$. Let $D \subseteq \Z^d$ be a finite domain. The function $H_{\theta |D}: \Sigma^{N_\theta(D)} \rightarrow \Sigma^D$ is a \emph{non-uniform \textsc{ca} over finite domain} $D$, which maps $c_{|N_\theta (D)}$, where $c \in \Sigma^\Z$, to
\begin{align*}
H_{\theta|D} (c_{|N_\theta (D)}) = H_\theta (c)_{|D}.
\end{align*}
\end{definition}

\begin{definition}
Let $H_\theta$ be a \textsc{nuca}. We say $H_\theta$ is \emph{balanced}, when for any finite domain $D\subseteq \Z$ and pattern $p \in \Sigma^D$, if $E = N_\theta(D)$ and $\mathcal{A} = \set{q \in \Sigma^E}{H_{\theta|D}(q) = p}$, then $|\mathcal{A}| = |\Sigma|^{|E|-|D|}$.
\end{definition}

%\begin{theorem}
%There exists a rule set $\R$ over state set $\Sigma$ and neighbourhood $N$, and rule distribution $\theta \in \R^\Z$, such that $H_\theta$ is surjective, but there also are finite domain $D \subseteq \Z$ and finite pattern $p' \in \Sigma^D$ such that if $t$ is the number of patterns $p \in N(D)$ with $H_{\theta|D} (p) = p'$, then $t \neq |\Sigma|^{|N(D)|- |D|}$.
%\end{theorem}

\begin{example}
Let $\Sigma = \Z_2$ and $\R = \{f,g\}$ be a set of radius $1$ rules, where
\begin{align*}
f(a,b,c) = a \oplus b \oplus c, \:
g(a,b,c) = \max (b,c), 
\end{align*}
Let $\theta \in \R^\Z$ be the rule distribution such that
\begin{align*}
\theta (x) &= \begin{cases}
g, \text{ if } x=0,\\
f, \text{ if } \text{otherwise.}
\end{cases}
\end{align*}
First, let's show $H_\theta$ is surjective. Let $c \in \Sigma^\Z$ be any configuration. Then let $e \in \Sigma^\Z$ be a configuration such that 
\begin{align*}
e(x) = \begin{cases}
0, \text{ if } x=-1 \text{ or } x=1\\
c(x), \text{ if } x=0,\\
c(x+1) \oplus e(x+1) \oplus e(x+2), \text{ if } x < -1,\\
c(x-1) \oplus e(x-1) \oplus e(x-2), \text{ if } x>1,
\end{cases}
\end{align*}
where $x \in \Z$. Now if $x \neq 0$ clearly $H_\theta (e)(x) = c(x)$, and if $x = 0$, $H_\theta (e)(x) = \max (c(x), 0) = c(x)$. Therefore $H_\theta (e) = c$, meaning $H_\theta$ is surjective.

On the other hand, let $D =\{0\}$ and $p' \in \Sigma^D$ be such that $p'(0) = 1$. Now the patterns $p \in \Sigma^{N_\theta(D)}$ such that $H_{\theta|D} (p) = p'$ are $001,010,011,101,110$ and $111$, for a total of $t = 6$ patterns. Then clearly $t = 6 \neq 4 = 2^2 =  |\Sigma|^{|N(D)|-|D|}$. 
\end{example}

Note that the \textsc{nuca} in the previous example is also pre-injective, that is, all configurations which differ in only finitely many cells have different images. Unlike in uniform \textsc{ca}, pre-injectivity and surjectivity aren't necessarily equivalent in \textsc{nuca} \cite{goe}.

It can be shown that a surjective \textsc{nuca} of any dimension with a uniformly recurrent rule distribution has the balance property, with a proof very similar to the proof on uniform \textsc{ca}. In the one-dimensional case, it is enough for the rule distribution to be recurrent. Note that surjectivity and pre-injectivity are equivalent in \textsc{nuca} with recurrent rule distributions \cite{goe}.

%TODO move this to appendix
%\begin{lemma} \label{ineq1} 
%Let $d,s,n,r \in \Z_+$. For all sufficiently large $k \in \Z_+$, it holds that
%\begin{align*}
%(s^{n^d}-1)^{k^d} < s^{(kn-2r)^d}.
%\end{align*}
%\end{lemma} 
%
%\begin{theorem} \label{balance1}
%Let $\theta \in \R^{\Z^d}$ be a uniformly recurrent rule distribution such that $H_\theta$ is surjective. Let $D\subseteq \Z^d$ be a finite domain. For every pattern $p' \in \Sigma^D$ the number of patterns $p \in \Sigma^{N(D)}$ such that $H_{\theta|D}(p) = p'$ is $s^{|N(D)|-|D|}$, where $s = |\Sigma|$.
%\end{theorem}
%
%\begin{proof}
%Without loss of generality, we can assume $N$ is a radius $r$ neighbourhood. Suppose there is $p' \in \Sigma^D$ such that the number of patterns $p \in \Sigma^{N(D)}$ that map to $p'$ is $t \neq s^{|N(D)|-|D|}$. Because $\theta$ is uniformly recurrent, there is $n \in \Z_+$ such that a copy of $\theta_{|D}$ appears in every width $n-2r$ hypercube in $\theta$.
%
%Let $k \in \Z_+$ be arbitrary. Consider finite domain $C \subseteq \Z^d$ that is a hypercube of size $(kn)^d$, partitioned into $k^d$ disjoint hypercubes of width $n$. Let $C'$ be the hypercube of size $(kn - 2r)^d$ centered inside $C$.
%
%Let $E'$ be a width $n-2r$ hypercube co-centric with one of the width $n$ hypercubes $E$. Then $\theta_{|E'}$ contains a copy of $\theta_{|D}$. This is illustrated in Figure \ref{fig_balance}. There are  $s^{|E'|-|D|}$ patterns in domain $E'$ that have $p'$ as a subpattern on top of a selected copy of $\theta_{|D}$. In total these have $t' = t\cdot s^{|E|-|N(D)|}$ pre-images in the domain $E$. Then if each pattern in the domain $E'$ would have $s^{|E|-|E'|}$ pre-images in $E$, we'd have
%\begin{align*}
%s^{|E'|-|D|}\cdot s^{|E|-|E'|} = t \cdot s^{|E|-|N(D)|},
%\end{align*}
%which would imply $t = s^{|N(D)|-|D|}$, a contradition. Therefore $t' \neq s^{|E|-|E'|}$.
%
%\begin{figure}
%\begin{center}
%\includegraphics[scale=1]{"balance.pdf"}
%\caption{The co-centric width $kn$ and $kn-2r$ hypercubes $C$ and $C'$, partitioned into $k^d$ hypercubes of width $n$. Copies of $\theta_{|D}$ (black squares) are found in width $n-2r$ hyper co-centric with the $n$ wide ones.} \label{fig_balance}
%\end{center}
%\end{figure}
%
%Additionally, without loss of generality, we can assume that $t' < s^{|E|-|E'|}$. This is because the number of patterns in domains $E$ and $E'$ are $s^{|E|}$ and $s^{|E'|}$ respectively, so if every pattern in domain $E'$ would have at least $s^{|E|-|E'|}$ pre-images, then every pattern in $E'$ would have exactly $s^{|E|-|E'|}$ pre-images, which would be a contradiction.
%
%Next, let $\mathcal{A}\subseteq \Sigma^{C'}$ be the set of patterns such that there's a copy of pattern $p$ on top of a copy of $\theta_{|D}$ in each of the $k^d$ width $n-2r$ hypercubes. Then $|\mathcal{A}| = s^{|C'|- k^d|D|}$. Let $\mathcal{B}\subseteq \Sigma^{C}$ be the set of patterns whose image under $H_{\theta|C'}$ is in $\mathcal{A}$. They are exactly the patterns that contain one of the $t'$ pre-images of the $p'$-bearing width $n-2r$ hypercube patterns, in each of the $k^d$ width $n$ hypercubes. Therefore $|\mathcal{B}|=(t')^{k^d}$.
%
%Now it remains to show that $|\mathcal{B}|<|\mathcal{A}|$ for sufficiently large $k$. We have $|E|= n^d$ and $|C'| = (kn-2r)^{k^d}$, so by Lemma \ref{ineq1}, $(s^{|E|}-1)<s^{|C'|}$. Then
%\begin{align*}
%|\mathcal{B}| &= (t')^{k^d} \leq (s^{|E|-|E'|}-1)^{k^d} \leq (s^{|E|-|E'|}-s^{-|E'|})^{k^d} \\
%&= s^{k^d|E'|}(s^{|E|}-1)^{k^d} < s^{|C'|-k^d|E'|} \leq s^{|C'|-k^d|D|}=|\mathcal{A}|.
%\end{align*}
%The last inequality is due to the fact that $|E'| \leq |D|$. Now because $|\mathcal{B}|<|\mathcal{A}|$, there's a pattern with no pre-image, meaning $H_\theta$ is not surjective. Then if $H_\theta$ is surjective, the theorem statement holds.
%\end{proof}


\begin{lemma} \label{ineq2} 
Let $s,n,r \in \Z_+$. For all sufficiently large $m \in \Z_+$, it holds that
\begin{align*}
(s^n-1)^m s^{k-nm} < s^{k-2r},
\end{align*}
for all $k \in \Z_+$. 
\end{lemma} 

\begin{theorem}
Let $\theta \in \mathcal{R}^\Z$ be a recurrent rule distribution such that $H_\theta$ is surjective. Then $H_\theta$ is balanced.
\end{theorem}

\begin{proof}
Let $D\subseteq \Z$ be a finite domain and $E=N_\theta (D)$. Without loss of generality, we can assume all rules in $\R$ have radius $r$. Suppose there is $p \in \Sigma^D$ such that the number of $H_{\theta|D}$ pre-images is $t\neq s^{|E|-|D|}$. We can select $p$ such that $t< s^{|E|-|D|}$, because if every pattern in domain $D$ were to have at least $s^{|E|-|D|}$ patterns, since there are only $s^E$ patterns in domain $E$, every pattern would have exactly $s^{|E|-|D|}$ pre-images, contradicting our assumption.

Because $\theta$ is recurrent, it contains infinitely many copies of $\theta_{|D}$. Let $n \in \Z_+$ be such that $|D|=n-2r$. Let then $m\in \Z_+$ be large enough that $(s^n-1)^m s^{k-nm}<s^{k-2r}$ for any $k\in \Z_+$. By Lemma \ref{ineq2} such a number exists.

Let $D'$ be a $k-2r$ wide segment where $k\in \Z_+$ is large enough that a $\theta_{|D'}$ contains $m$ disjoint copies of $\theta_{|D}$ that are at least $2r$ apart, and let $E'= N_\theta (D')$. Of course now $|E'|=k$. Let $\mathcal{B} \subseteq \Sigma^{D'}$ be the set of patterns such that there is a copy of $p$ on each disjoint copy of $\theta_{|D}$. Because the rest of the cells may be freely chosen, $|\mathcal{B}|=s^{|D'|-m|D|} = s^{k-2r-m(n-2r)}$.

Let then $\mathcal{A} = \set{q \in \Sigma^{E'}}{H_{\theta|D'}(q) \in \mathcal{B}}$. Then $|\mathcal{A}|=t^m \cdot s^{|E'|-m|E|} = t^m \cdot s^{k-mn}$, because on each $n$ wide segment centered the copies of $\theta_{|D}$ must be a copy of one of the $t$ pre-images of $p$, and the other cells may be chosen freely. Now it remains to show that $|\mathcal{A}| < |\mathcal{B}|$. Now because $t\leq s^{2r}-1$, using the previous, we have 
\begin{align*}
|\mathcal{A}| &= t^m \cdot s^{k-mn} \leq (s^{2r}-1)^m \cdot s^{k-mn}\leq (s^{2r}- s^{-(n-2r)})^m\cdot s^{k-mn}\\
&= (s^n - 1)^m  \cdot s^{k-mn} \cdot s^{-m(n-2r)} < s^{k-2r} \cdot s^{-m(n-2r)} = |\mathcal{B}|.
\end{align*}
Hence $H_\theta$ is not surjective. Then if $H_\theta$ is surjective, the theorem statement holds.
\qed
\end{proof}

Though surjectivity is not enough to guarantee balance, bijectivity is. In fact, it is enough to assume that every configuration has the same number of pre-images. This also holds for any dimension, though we do not show it here.

\begin{theorem} \label{bij_bal}
Let $\theta \in \mathcal{R}^\Z$ be a rule distribution such that $H_\theta$ is bijective. Then $H_\theta$ is balanced.
\end{theorem}

\begin{proof}
Let $D\subseteq \Z$ be a finite segment and $p \in \Sigma^D$ a pattern. Let $E = N_\theta(D)$ and $\mathcal{A}_p = \set{q \in \Sigma^E}{H_{\theta|D}(q) = p}$.
Without loss of generality, we can assume $D$ is a contiguous segment. Let $D' = \Z \setminus N_\theta(E)$ and $E' = \Z \setminus E$. Let $C = \Z \setminus (D \cup D')$. Let $p' \in \Sigma^{D'}$ be a fixed (infinite) pattern and let $\mathcal{B}_{p'} = \set{q' \in \Sigma^{E'}}{H_{\theta|D'} (q') = p'}$. Let $p \in \Sigma^D$ be a finite pattern. Then consider $\mathcal{C}_{p,p'} = \set{c \in \Sigma^\Z}{c_{|D} = p, c_{|D'} = p' }$ and $\mathcal{C}_{p,p'}' = \set{c' \in \Sigma^\Z}{H_\theta(c') \in \mathcal{C}_{p,p'}}$. Obviously $|\mathcal{C}_{p,p'}| = |\Sigma|^{|C|}$. On the other hand, we have $|\mathcal{C}_{p,p'}'|=|\mathcal{A}_p|\cdot |\mathcal{B}_{p'}|$. Finally, because $H_\theta$ is bijective, we have $|\mathcal{C}_{p,p'}| = |\mathcal{C}_{p,p'}'|$. Therefore
\begin{align*}
|\mathcal{A}_p|\cdot |\mathcal{B}_{p'}| = |\mathcal{C}_{p,p'}'| = |\mathcal{C}_{p,p'}| = |\Sigma|^{|C|},
\end{align*}
which gives $|\mathcal{A}_p| = \frac{|\Sigma|^{|C|}}{|\mathcal{B}_{p'}|}$. Choice of $p'$ and $C$ was independent of the choice of $p$, and therefore this is true for any $p \in \Sigma^D$, and therefore every pattern in the domain $D$ has the same number of pre-images. Hence we have  $|\mathcal{A}_p| = |\Sigma|^{|E|-|D|}$.
\qed
\end{proof}


\begin{figure}
\begin{center}
	\includegraphics[scale=0.9]{bij_bal.pdf}
	\caption{Naming scheme of the domains used in the proof of Theorem \ref{bij_bal}.} \label{fig_bij_bal}
\end{center}
\end{figure}


\section{Sensitivity and equicontinuity}

Sensitivity and equicontinuity are properties of general dynamical systems. To define them, some topological concepts are needed. 

\begin{definition}
Let $\Sigma$ be a state set, $c \in \Sigma^\Z$ a configuration and $D \subseteq \Sigma$ a finite domain. We denote a \emph{cylinder} as
\begin{align*}
\mathrm{Cyl} (c,D) = \set{e \in \Sigma^\Z}{c_{|D} = e_{|D}},
\end{align*}
when $\Sigma$ is clear from context.
\end{definition}

It's well known that cylinders form a topological basis for the Cantor space \cite{kari}, that is, every open set is a union of cylinders. Topologically, a set $A\subseteq \Sigma^\Z$ is dense if every cylinder contains an element of $A$, and residual if it is a countable intersection of dense open sets. Because the Cantor space is a Baire space, all residual sets are also dense. %TODO add cite

The following definitions are equivalent with the topological definitions of equicontinuity and sensitivity. They can naturally be extended to any dimension, but here we focus on only 1-dimensional \textsc{nuca}.

\begin{definition}
Let $H_\theta$ be a \textsc{nuca} with the state set $\Sigma$. The configuration $c \in \Sigma^\Z$ is an \emph{equicontinuity point} if for every finite $E \subseteq \Z$ there is a finite $D\subseteq \Z^d$ such that
\begin{align*}
\forall  e \in \mathrm{Cyl}(c,D)  , n\in \Z_+: H_\theta^n (e) \in \mathrm{Cyl}(H_\theta^n(c),E).
\end{align*}
We denote the set of equicontinuity points as $\mathcal{E}$. We say $H_\theta$ is \emph{equicontinuous}, if $\mathcal{E} = \Sigma^\Z$. If $\mathcal{E}$ is residual, we say $H_\theta$ is \emph{almost equicontinuous}.
\end{definition}

In uniform \textsc{ca} we find that only eventually temporally periodic \textsc{ca} are equicontinuous. However, there are (reversible) equicontinuous \textsc{nuca} which have no periodic configurations \cite{kamilya}.

\begin{definition}
Let $H_\theta$ be a \textsc{nuca} with the state set $\Sigma$. The \textsc{nuca} is called \emph{sensitive} if there exists a finite $E\subseteq \Z$ such that for every $c \in \Sigma^\Z$ and every finite $D \subseteq \Z$,
\begin{align*}
\exists  e\in \mathrm{Cyl} (c,D),n \in \Z_+ :  H_\theta^n (e) \notin \mathrm{Cyl}(H_\theta^n(c),E).
\end{align*}
\end{definition}

In 1-dimensional uniform \textsc{ca}, we find that either $\mathcal{E} = \emptyset$, which is equivalent with sensitivity, or $\mathcal{E}$ is residual. 
However, there are 2-dimensional uniform \textsc{ca} whose set of equicontinuity points is neither empty nor residual, as well as ones which have no equicontinuity points but are not sensitive. \cite{equisense}

In the 1-dimensional non-uniform case, the dichotomy found in uniform \textsc{ca} does not hold, and in fact there exist both 1D \textsc{nuca} whose set of equicontinuity points is neither empty nor residual, and 1D \textsc{nuca} which have no equicontinuity points but are not sensitive. 

\begin{theorem} \label{equicounter}
There is a rule set $\R$ and rule distribution $\theta \in \R^\Z$ such that if $\mathcal{E}$ is the set of equicontinuity points of $H_\theta$, then $\mathcal{E} \neq \emptyset$ and $\mathcal{E}$ is not residual.
\end{theorem}

\begin{proof}
Let $\Sigma = \{0,1\}$ and $\R = \{\mathrm{id}, \tau\}$ where $\mathrm{id}$ is the identity rule and and $\tau$ is the leftward traffic rule, a radius $1$ which maps
\begin{align*}
\tau (a,b,c) = \begin{cases}
0, \text{ if } a=0, b=1, \\
1, \text{ if } b=0, c=1, \\
b, \text{ otherwise,}
\end{cases}
\end{align*}
for any $a,b,c\in\Sigma$. Let $\theta \in \R^\Z$ be such that $\theta(x) = \tau$ if $x> 0$ and $\theta(x)=\mathrm{id}$ if $x\leq 0$. Let $\mathcal{E}$ be the set of equicontinuity points of $H_\theta$. The operation of $H_\theta$ is illustrated in Figure \ref{fig_equi1}.

\begin{figure}
\begin{center}
\includegraphics[scale=0.7]{"traffic.pdf"}
\caption{Space-time diagram of the \textsc{nuca} defined in the proof of Theorem \ref{equicounter}. On the left, cell $0$ is set to "1" and doesn't let traffic pass, whereas on the right its set to "0" and annihilates traffic.} \label{fig_equi1}
\end{center}
\end{figure}

%The informal explanation of the proof's idea is that cell 0 acts as a fixed "gate" for the leftward traffic. If cell 0 is set to "1" it blocks traffic, and if it's set to "0" it annihilates traffic. When traffic is blocked, any traffic accumulated at the origin is frozen. This means a configuration with "1" in all non-negative cells an equicontinuity point, because any observation window containing the origin is frozen. On the other hand, any configuration with "0" at cell 0 is not an equicontinuity point, because traffic from arbitrarily far can reach the origin. The function of this \textsc{nuca} is illustrated in Figure \ref{fig_equi1}\\

First to show that $\mathcal{E} \neq \emptyset$. Let $c_1 \in \Sigma^\Z$ be a configuration of all 1's. Because every cell in $c_1$ sees only 1's, clearly $H_\theta (c_1) = c_1$. Let $E \subseteq \Z$ be a finite segment and let $k\in \Z$ be the rightmost cell of $E$ and let $D = E \cup [0,k]$. Let then $e \in \mathrm{Cyl}(c_1,D)$. For all $x\leq 0$ obviously $H_\theta^n (e) (x) = e(x)$ for any $n \geq 1$. Then if $k\leq 0$, we have $H_\theta^n (e) (x) = e(x) = c(x) = H_\theta^n (c)(x) $ for all $x \in E$. 

Suppose then $k> 0$. For any cell $x>0$, if $e(x)=1$, then $H_\theta(e)(x)=0$ only if $e(x-1) = 0$. However because for all $x \in [0,k]$, $e(x)=1$, and by the previous $H_\theta(e)(0)=1$, then also $H_\theta(e)(x) = 1$. Now we have $H_\theta (e) \in \mathrm{Cyl} (c_1, D)$ and because $e$ was arbitrary, then $H_\theta (\mathrm{Cyl} (c_1, D) ) \subseteq \mathrm{Cyl} (c_1, D)$ and hence $H_\theta ^n (e) \in \mathrm{Cyl} (c_1, D) \subseteq \mathrm{Cyl} (c_1, E)$ for all $n>0$. Therefore $c_1 \in \mathcal{E}$.\\

Now to show that $\mathcal{E}$ is not residual. Because residual sets in a Baire space are dense, it is enough to show that $\mathcal{E}$ is not dense. Let $c \in \Sigma^\Z$ be any configuration such that $c(0) = 0$. We show that $c \notin \mathcal{E}$. Let $E = \{1\}$ and $D \subseteq \Z$ be any finite domain. Without loss fo generality we can assume $1 \in D$. Informally, the idea is that the leftward traffic is annihilated at cell $0$ and traffic from arbitrarily far to the left can then arrive at $E$. 

Let $e_0,e_1\in \mathrm{Cyl}(c,D)$ be such that for all $x \in \Z \setminus D$, $e_0(x) = 0$ and $e_1(x) = 1$. Firstly, there are finitely many $x>0$ such that $e_0(x) = 1$. Also the number of such cells can only decrease, since under the rule $\tau$, when a "0" cell turns into a "1" cell, another "1" cell turns into a "0" cell. Then for any $t\in \Z_+$ such that $H_\theta^t (e_0)(1) = 1$, the number of cells $x>0$ such that $H_\theta^{t+1} (e_0)(x) = 1$ is greater than the number of such cells with $H_\theta^{t} (e_0)(x) = 1$. On the other hand if for some $k>0$ and $H_\theta^t (e_0)(i) = 0$ for all $i \in [1,k]$ and $H_\theta^t (e_0)(i) = 1$ and $H_\theta^t (e_0)(k+1)=1$, then $H_\theta^{t+k} (e_0)(1) = 1$. Hence there is some $n\geq 0$ such that for all $t\geq n$, $H_\theta^{t} (e_0)(1) = 0$, because the number of positive "1" cells decreases to 0.

On the other hand, there are infinitely many $t \geq 0$ such that $H_\theta^t (e_1)(1) = 1$, because there are infinitely many "1" cells to the right of 1.   Now there must be some $t\geq 0$ such that $H_\theta^t(e_0)(1) \neq H_\theta^t(e_1)(1)$. Therefore either $H_\theta^t (e_0) \notin \mathrm{Cyl}(c,E)$ or $H_\theta^t(e_1) \notin \mathrm{Cyl}(c,E)$. In either case, because $D$ was arbitrary, this means $c$ is not an equicontinuity point. Now since no configuration with "0" at cell 0 is an equicontinuity point, $\mathcal{E}$ is not dense, and therefore not residual.

%Under $\tau$, a "1" cell becomes a "0" cell if it sees a "0" to its left. Then because for all $n>0$ we have $H_\theta ^n (e)(0) = 0$, by induction for every $x>0$ there's infinitely many $t>0$ such that $H_\theta^t (e)(x)=0$. If there is a cell $x>0$ such that for all $y\geq x$, $e(y)=0$, then clearly for all $t>0$, $H_\theta ^t (e)(x) = 0$. Then inductively there is such $n$ for every $x>0$.

%Now, let $e_0,e_1\in \mathrm{Cyl}(c,D)$ be such that for all $x \in \Z \setminus D$, $e_0(x) = 0$ and $e_1(x) = 1$. By the previous, there is $n\geq 0$ such that for all $t\geq n$, $H_\theta ^t (e_0)(1) = 0$. On the other hand, for any $t\geq 0$,  there are infinitely many $x>0$ such that $H_\theta ^t (e_1)(x) =1$. Then if for some $t\geq 0$ there's $k> 0$ such that all for $x \in [1,k]$, $H_\theta ^t (e_1)(x) = 0$ and $H_\theta^t (e_1)(k+1) = 1$, then clearly $H_\theta ^{t+k} (e_1)(1) =1$. Therefore there's infinitely many $t\geq 0$ such that $H_\theta^t(e_1)(1) = 1$.


\qed
\end{proof}



\begin{theorem} \label{equicounter2} 
There is a rule set $\mathcal{R}$ and rule distribution $\theta \in \mathcal{\Z}$ such that $H_\theta$ is not sensitive, and if $\mathcal{E}$ is the set of equicontinuity points of $H_\theta$, then $\mathcal{E} = \emptyset$.
\end{theorem}

\begin{proof}
Let $\Sigma = \{0,1,2,3 \}$ and $\mathcal{R} = \{\id, f \}$ a rule set, where $\id$ is the identity rule, and $f:\Sigma^3 \rightarrow \Sigma$ is a radius-1 rule that maps
\begin{align*}
&(a,0,3)\mapsto 3, (1,1,1)\mapsto 1, (1,1,2)\mapsto 1, (a,1,3)\mapsto 3,\\
&(1,2,c)\mapsto 2, (a,3,3)\mapsto 3, \text{ and } (a,b,c)\mapsto 0 \text{ otherwise},
\end{align*} 
where $a,b,c\in \Sigma$. Let $\theta \in \mathcal{R}^\Z$ be such that $\theta (x) = \mathrm{id}$ if $x\leq 0$ and $\theta (x) = f$ if $x>0$. The operation of this \textsc{nuca} is illustrated in Figure \ref{fig_equi2}.

\begin{figure}
\begin{center}
\includegraphics[scale=0.8]{"equi2"}
\caption{Space-time diagrams of the \textsc{nuca} defined in the proof of Theorem \ref{equicounter2}. The left diagram illustrates how a segment arbitrarily far from the origin can be "protected" from 3-cells with a blocking pattern, and the right diagram illustrates why said blocking pattern can't be copied to protect every cell from 3-cells.} \label{fig_equi2}
\end{center}
\end{figure}


First to show that $H_\theta$ is not sensitive: let $E\subseteq [i,j] \subseteq \Z$ be an arbitrary finite domain. Let $c\in \Sigma^\Z$ be such that $c(j) = 2$ and $c(x) = 1$ if $x\neq j$, and let then $D=[\min (i,0),j]$. Let $e \in \cyl (c,D)$. If $x\in[i, 0]$, then obviously $H_\theta (e)(x) = c(x)$, because $\theta (x) = \id$. For every $x \in [1,j-1]$, $H_\theta (e)(x) = f(1,1,a) = 1 = c(x)$, where $a \in \{1,2\}$ and $H_\theta (e) (j) = f(1,2,b) = 2 = c(j)$, where $b \in \Sigma$. Then $H_\theta (e) \in \cyl (c,D)$, and because $e$ was an arbitrary element of $\cyl (c,D)$, then $H_\theta ^n (e) \in \cyl (c,D)$ for all $n>0$. Since also $c \in \cyl (c,D)$, and $E \subseteq D$, we have $H_\theta ^n (e) \in \cyl ( H_\theta^n (c),D )\subseteq \cyl (H_\theta ^n (c), E)$ for all $n>0$. Therefore $H_\theta$ is not sensitive.

Then to show that $\mathcal{E} = \emptyset$. First we show that if a "2"-cell (on the right side) has any non-"1"-cell between it and the origin, it will eventually become a "0"-cell. %this is true but it's enough that its not a 2-cell. then you don't have to also show that 3's can't penetrate the 2 again
Let $c \in \Sigma^\Z$ be a configuration with  $x_1,x_2 \in \Z$, $0\leq x_1< x_2$ such that $c(x_1)\neq 1$, $c(x_2)=2$ and $c(x) = 1$ when $x_1 < x < x_2$. Suppose for some $n \in [0, x_2-x_1-2]$ we have 
\begin{align*}
H_\theta ^n (c)(x_1+n) \neq 1, \:
H_\theta ^n (c)(x_2) = 2, \:
H_\theta ^n (c)(x) = 1,
\end{align*}
if $x_1+n < x < x_2$. Then under the local rule $f$,
\begin{align*}
H_\theta ^{n+1} (c)(x_1+n+1) = 0,\:
H_\theta ^{x+1} (c)(x_2) = 2, \:
H_\theta ^{n+1} (c)(x) = 1,
\end{align*}
if $x_1+n+1 < x < x_2$. Then by induction, $H_\theta ^{x_2-x_1-1}(c)(x_2-1)=0$ and $H_\theta ^{x_2-x_1-1}(c)(x_2)=2$, and hence $H_\theta ^{x_2-x_1}(c)(x_2)=0$.

Let now $c \in \Sigma^\Z$. We show that $c \notin \mathcal{E}$. Let $x\geq 1$ be the smallest integer such that $c(x-1)=2$. If no such cell exists, we let $x=0$. Let $E=\{x\}$, and let $D=[y,z]$ with any $y,z \in \Z$, $y\leq z$. Without loss of generality, we can assume $x \in D$. Let $e_1,e_2 \in \cyl (c,D)$ be such that if $k \notin D$, then $e_1(k) = 0$ and $e_2(k) = 3$.

By the previous, because no non-"2"-cell can become a "2"-cell, there is $n\geq 0$ such that for all $k \in [x,z]$, $H_\theta ^n (e_1)(k) \neq 2$ and $H_\theta ^n (e_2)(k) \neq 2$. Then there are only finitely many "3"-cells to the right of $x$ in $H_\theta ^n (e_1)$, so for some $n_1 \geq n$, for all $t\geq n_1$ we have $H_\theta ^t (e_1)(x) \neq 3$. On the other hand, because there are infinitely many "3"-cells to the right of $x$ in $H_\theta ^n (e_2)$, for some $n_2 \geq n$, for infinitely many $t > n_2$, $H_\theta ^t (e_2)(x) = 3$. Then there is some $t>0$ such that $H_\theta ^t (e_1)(x) \neq 3 = H_\theta ^t (e_2)(x)$, and therefore either $H_\theta ^t (e_1) \notin \cyl (H_\theta ^t (c),E)$ or $H_\theta ^t (e_2) \notin \cyl (H_\theta ^t (c),E)$. Since this is true for any finite segment $D$, $c \notin \mathcal{E}$, and therefore $\mathcal{E} = \emptyset$.
\qed
\end{proof}


It is still open whether there exist counter-examples that have uniformly recurrent rule distributions.

\section{Expansivity}


In \textsc{ca}, expansivity is a much stronger notion of sensitivity. A \textsc{(nu)ca} is expansive if there's a fixed window where any difference between configurations propagates.

\begin{definition}
Let $\theta \in \mathcal{R}^\Z$ be a rule distribution. The \textsc{nuca} $H_\theta$ is \emph{positively expansive} if there's some finite domain $E \subseteq \Z$ such that for any $c,e\in \Sigma^\Z$, $c \neq e$, there is $n\in \Z_+$ such that $H_\theta ^n (e) \notin \cyl (H_\theta ^n (c),E)$.

The \textsc{nuca} $H_\theta$ is \emph{expansive} if it's bijective and there's a finite domain $E \subseteq \Z$ such that for any $c,e\in \Sigma^\Z$, $c \neq e$, there is $n\in \Z$ such that $H_\theta ^n (e) \notin \cyl (H_\theta ^n (c),E)$.
\end{definition}

In uniform \textsc{ca}, the usual proof that expansivite \textsc{ca} are sensitive utilizes the denseness of periodic points. In bijective \textsc{nuca}, periodic points are not necessarily dense \cite{kamilya}, so we use the Poincaré recurrence theorem instead. Let's first recall some definitions from ergodic theory.

\begin{definition}
Let $\Sigma$ be a state set. A Borel measure $\mu: \mathcal{B} \rightarrow \mathbb{R}$, where $\mathcal{B}$ is a Borel algebra over $\Sigma^\Z$, is \emph{uniform} if it maps 
\begin{align*}
\cyl (c,D) \mapsto \left( \frac{1}{|\Sigma|} \right) ^{|D|},
\end{align*}
where $c \in \Sigma^\Z$ and $D\subseteq \Z$ is a finite domain.\\ %TODO add cite

Let $H_\theta:\Sigma^\Z \rightarrow \Sigma^\Z$ be a \textsc{nuca}. We denote $H_\theta(\mu):\mathcal{C}\rightarrow \mathbb{R}$ as the measure that gives $H_\theta(\mu)(A) = \mu (H_\theta^{-1} (A))$ for any $A \in \mathcal{B}$. We say that measure $\mu$ is \emph{invariant under} $H_\theta$ if $H_\theta(\mu) = \mu$.
\end{definition}

It's well known that the value of a measure on cylinders uniquely determines it, that is, we need to only define a measure on cylinders. \cite{ergodic}
The following proof is essentially identical to the proof on uniform automata.
 
\begin{lemma}
Let $H_\theta:\Sigma^\Z \rightarrow \Sigma^\Z$ be a bijective \textsc{nuca}. The uniform probability measure $\mu$ is invariant under $H_\theta$.
\end{lemma}

\begin{proof}
Let $C = \cyl (c, D)$ for some finite domain $D \subseteq \Z$ and $c \in \Sigma^\Z$. By Theorem \ref{bij_bal},
\begin{align*}
H_\theta^{-1} (C) = C_1 \cup C_2 \cup \ldots \cup C_n,
\end{align*}
where $C_i$ are disjoint cylinders with some finite domain $E \subseteq \Z$, and $n=|\Sigma|^{|E|-|D|}$. Then due to the additivity of measures, \cite{ergodic}
we have
\begin{align*}
\mu (H_\theta^{-1} (C) ) = \mu (C_1) + \mu (C_2) + \ldots + \mu(C_n) = n \left( \frac{1}{|\Sigma|} \right)^{|E|} = \left( \frac{1}{|\Sigma|} \right)^{|D|} = \mu (C).
\end{align*}
Therefore $H_\theta (\mu) = \mu$.
\qed
\end{proof}

It is well know that if a dynamical system has an invariant measure of full support, recurrent points in it are dense. Therefore we state the following lemma without proof. \cite{ergodic}

\begin{definition}
Let $H_\theta:\Sigma^\Z \rightarrow \Sigma^\Z$ be a \textsc{nuca}. A configuration $c \in \Sigma^\Z$ is \emph{temporally recurrent} if for any $D \subseteq \Z$ there exists $n\in \Z_+$ such that $H_\theta^n (c) \in \cyl (c,D)$.
\end{definition}

\begin{lemma}[Poincaré]\label{poincare}
Let $H_\theta:\Sigma^\Z \rightarrow \Sigma^\Z$ be a bijective \textsc{nuca}. Then temporally recurrent configurations of $H_\theta$ are dense, that is, for any cylinder $C$ there is a configuration $c \in C$ that is temporally recurrent under $H_\theta$. 
\end{lemma}

Now we are ready to show that expansive \textsc{nuca} are indeed sensitive.

\begin{theorem}
Let $\theta \in \mathcal{R}^\Z$ be a rule distribution such that $H_\theta$ is expansive or positively expansive. Then $H_\theta$ is sensitive.
\end{theorem}

\begin{proof}
If $H_\theta$ is positively expansive, it is clearly sensitive by definition. Suppose then that (bijective) $H_\theta$ is expansive. Let $E \subseteq \Z$ be the observation window confirming expansivity, that is, a finite domain such that for any $c_1,c_2\in \Sigma^\Z$, $c_1 \neq c_2$, there is $n\in \Z$ such that $H_\theta ^n (c_2) \notin \cyl (H_\theta ^n (c_1),E)$. Let $D\subseteq \Z$ be a finite domain and $c_1 \in \Sigma^\Z$ a configuration. It is enough to show that for some $c_2 \in \cyl (c_1,D)$, $m \in \Z_+$ and $x \in E$, we have $H_\theta^m (c_1)(x) \neq H_\theta^m (c_2)(x)$.

By expansivity there is $n\in \Z$ such that $H_\theta^n (c_1)(x) \neq H_\theta^n (c_2)(x)$ for some $x\in E$. If $n\geq 0$ the proof is complete, so assume $n<0$. Consider the space $(\Sigma \times \Sigma)^\Z$ and \textsc{nuca} $H_\phi:(\Sigma \times \Sigma)^\Z \rightarrow (\Sigma \times \Sigma)^\Z$ such that for $e \in (\Sigma \times \Sigma)^\Z$ and $e_1,e_2 \in \Sigma^\Z$ where $e(y) = (e_1(y),e_2(y))$ for all $y\in \Z$, $H_\phi$ maps $e$ such that
\begin{align*}
H_\phi(e)(y) = (H_\theta (e_1)(y),H_\theta (e_2)(y)),
\end{align*}
for all $y\in \Z$. Obviously such $e_1,e_2$ exist for any $e$, and $H_\phi$ is a \textsc{nuca}. 

Because $H_\theta$ is bijective, clearly $H_\phi$ is too. Then by Lemma \ref{poincare}, temporally recurrent configurations of $H_\phi$ are dense. Then there's a temporally recurrent $e \in (\Sigma \times \Sigma)^\Z$ with $e_1,e_2 \in \cyl (c, D)$ such that $e(y) = (e_1(y),e_2(y))$ for all $y \in \Z$ and $H_\theta^n (e_1)(x) = H_\theta^n (c_1)(x)$, $H_\theta^n (e_2)(x) = H_\theta^n (c_2)(x)$ by the continuity of $H_\theta$. Because $e$ is temporally recurrent, there is some $m>0$ such that $H_\theta^m (e_1)(x) = H_\theta^n (e_1)(x)$ and $H_\theta^m (e_2)(x) = H_\theta^n (e_2)(x)$. Then because 
\begin{align*}
H_\theta^m (e_1)(x) = H_\theta^n (c_1)(x) \neq H_\theta^n (c_2)(x) = H_\theta^m (e_2)(x),
\end{align*}
either $H_\theta^m (e_1)(x) \neq H_\theta^m (c_1)(x)$ or $H_\theta^m (e_2)(x) \neq H_\theta^m (c_1)(x)$. Therefore $H_\theta$ is sensitive.
\end{proof}

\section{Conclusions}

In summary, we've shown that a bijective \textsc{nuca} with a uniformly recurrent rule distribution is reversible and that this is not true of bijective \textsc{nuca} in general, even with recurrent rule distributions. We've shown that a bijective \textsc{nuca} is balanced, and that a surjective \textsc{nuca} with a recurrent a rule distribution is balanced, once again not true of all surjective \textsc{nuca}. 

We've also presented counter-examples for the usual dichotomy of sensitivity and equicontinuity in 1D \textsc{ca}, that is, a 1D \textsc{nuca} which is not sensitive and has no equicontinuity points, and one which has a non-residual and non-empty set of equicontinuity points. The question of whether there are such counter-examples with (uniformly) recurrent rule distributions remains open.  Finally, we've shown that (positively) expansive \textsc{nuca} are sensitive.

\subsubsection{Acknowledgements} 

We would like to acknowledge funding from the Magnus Ehrnrooth foundation and the Research Council of Finland, grant 354965.

%
% ---- Bibliography ----
%
% BibTeX users should specify bibliography style 'splncs04'.
% References will then be sorted and formatted in the correct style.
%
% \bibliographystyle{splncs04}
% \bibliography{mybibliography}
%
\begin{thebibliography}{8}


\bibitem{cecc} Cecchirini-Silberstein, T., Coornaert, M.: Cellular Automata and Groups. Springer Berlin, Heidelberg (2010)

\bibitem{dennuz}
Dennunzio, A., Formenti, E., Provillard, J.: Non-uniform Cellular Automata: Classes, dynamics, and decidability. Information and Computation \textbf{215}, 32--46 (2012)

\bibitem{goe} Paturi, P. (2023). On the Surjunctivity and the Garden of Eden Theorem for Non-uniform Cellular Automata. In: Manzoni, L., Mariot, L., Roy Chowdhury, D. (eds) Cellular Automata and Discrete Complex Systems. AUTOMATA 2023. Lecture Notes in Computer Science, vol 14152. Springer, Cham. \doi{10.1007/978-3-031-42250-8_2}

\bibitem{kari} Kari, K.: Theory of cellular automata: A survey. Theoretical Computer Science \textbf{334}(1--3), 3--33 (2005)

\bibitem{kamilya} Kamilya, S., Kari, J. Nilpotency and periodic points in non-uniform cellular automata. Acta Informatica \textbf{58}, 319--333 (2021). \doi{10.1007/s00236-020-00390-7}

\bibitem{equisense} Sablik, M., Theyssier, G. Topological Dynamics of Cellular Automata: Dimension Matters. Theory Comput Syst \textbf{48}, 693--714 (2011). \doi{10.1007/s00224-010-9255-x}

\bibitem{ergodic} Choe, G.: Computational Ergodic Theory. 1st edn. Springer Berlin, Heidelberg (2005) %TODO switch out this cite for better textbook

\bibitem{phung}  Phung, X.: On invertible and stably reversible non-uniform cellular automata, Theoretical Computer Science, \textbf{940}, Part B, 43--59 (2023)

%\bibitem{ref_article1}
%Author, F.: Article title. Journal \textbf{2}(5), 99--110 (2016)
%
%\bibitem{ref_lncs1}
%Author, F., Author, S.: Title of a proceedings paper. In: Editor,
%F., Editor, S. (eds.) CONFERENCE 2016, LNCS, vol. 9999, pp. 1--13.
%Springer, Heidelberg (2016). \doi{10.10007/1234567890}
%
%\bibitem{ref_book1}
%Author, F., Author, S., Author, T.: Book title. 2nd edn. Publisher,
%Location (1999)
%
%\bibitem{ref_proc1}
%Author, A.-B.: Contribution title. In: 9th International Proceedings
%on Proceedings, pp. 1--2. Publisher, Location (2010)
%
%\bibitem{ref_url1}
%LNCS Homepage, \url{http://www.springer.com/lncs}. Last accessed 4
%Oct 2017
\end{thebibliography}
\end{document}
